% Preamble for the Chinese two-column drafts (自动化学报 / Chinese conference style).
% Requires: \documentclass[twocolumn,10pt]{ctexart} before % Preamble for the Chinese two-column drafts (自动化学报 / Chinese conference style).
% Requires: \documentclass[twocolumn,10pt]{ctexart} before % Preamble for the Chinese two-column drafts (自动化学报 / Chinese conference style).
% Requires: \documentclass[twocolumn,10pt]{ctexart} before % Preamble for the Chinese two-column drafts (自动化学报 / Chinese conference style).
% Requires: \documentclass[twocolumn,10pt]{ctexart} before \input{preamble}.
\usepackage{amsmath,amssymb,amsthm,mathtools}
\usepackage{graphicx}
\usepackage{booktabs,multirow,array,tabularx}
\usepackage[caption=false]{subfig}
\usepackage{algorithm}
\usepackage{algpseudocode}
\usepackage[margin=2cm]{geometry}
\usepackage{url}
\usepackage{textcomp}
\usepackage{microtype}
\usepackage{xcolor}
\usepackage{bm}
\usepackage[hidelinks]{hyperref}

\setCJKmainfont{Noto Serif CJK SC}
\setCJKsansfont{Noto Sans CJK SC}
\setCJKmonofont{Noto Sans CJK SC}

\graphicspath{{figures/}}
\linespread{1.15}
\raggedbottom
\setlength{\emergencystretch}{2em}
\setlength{\columnsep}{0.75cm}
\renewcommand{\figurename}{图}
\renewcommand{\tablename}{表}
\renewcommand{\refname}{参考文献}
\floatname{algorithm}{算法}
\algrenewcommand\algorithmicrequire{\textbf{输入：}}
\algrenewcommand\algorithmicensure{\textbf{输出：}}

% Independent counters for lemmas / theorems / propositions so that the numbering
% matches the English source (Lemma 1, Theorem 1, Lemmas 2--5, Theorems 2--3,
% Propositions 1--3), see en/preamble.tex.
\newtheorem{theorem}{定理}
\newtheorem{lemma}{引理}
\newtheorem{proposition}{命题}
\newtheorem{corollary}{推论}
\theoremstyle{definition}
\newtheorem{definition}{定义}
\newtheorem{assumption}{假设}
\newtheorem{example}{例}
\theoremstyle{remark}
\newtheorem{remark}[theorem]{注}

\newcommand{\R}{\mathbb{R}}
\newcommand{\Sm}{\mathbb{S}}
\newcommand{\bq}{\mathbf{q}}
\newcommand{\bt}{\mathbf{t}}
\newcommand{\bx}{\mathbf{x}}
\newcommand{\bu}{\mathbf{u}}
\newcommand{\br}{\mathbf{r}}
\newcommand{\bM}{\mathbf{M}}
\newcommand{\bD}{\mathbf{D}}
\newcommand{\bG}{\mathbf{G}}
\newcommand{\bH}{\mathbf{H}}
\newcommand{\bE}{\mathbf{E}}
\newcommand{\bJ}{\mathbf{J}}
\newcommand{\vtwo}{v_2}
\newcommand{\tr}{\operatorname{tr}}
\newcommand{\rank}{\operatorname{rank}}
\newcommand{\diag}{\operatorname{diag}}
\newcommand{\Proj}{\operatorname{Proj}}
\newcommand{\argmin}{\operatorname*{arg\,min}}

\newcommand{\pend}[1]{\textcolor{red}{[#1]}}
\newcommand{\ev}[1]{\emph{#1}}
.
\usepackage{amsmath,amssymb,amsthm,mathtools}
\usepackage{graphicx}
\usepackage{booktabs,multirow,array,tabularx}
\usepackage[caption=false]{subfig}
\usepackage{algorithm}
\usepackage{algpseudocode}
\usepackage[margin=2cm]{geometry}
\usepackage{url}
\usepackage{textcomp}
\usepackage{microtype}
\usepackage{xcolor}
\usepackage{bm}
\usepackage[hidelinks]{hyperref}

\setCJKmainfont{Noto Serif CJK SC}
\setCJKsansfont{Noto Sans CJK SC}
\setCJKmonofont{Noto Sans CJK SC}

\graphicspath{{figures/}}
\linespread{1.15}
\raggedbottom
\setlength{\emergencystretch}{2em}
\setlength{\columnsep}{0.75cm}
\renewcommand{\figurename}{图}
\renewcommand{\tablename}{表}
\renewcommand{\refname}{参考文献}
\floatname{algorithm}{算法}
\algrenewcommand\algorithmicrequire{\textbf{输入：}}
\algrenewcommand\algorithmicensure{\textbf{输出：}}

% Independent counters for lemmas / theorems / propositions so that the numbering
% matches the English source (Lemma 1, Theorem 1, Lemmas 2--5, Theorems 2--3,
% Propositions 1--3), see en/preamble.tex.
\newtheorem{theorem}{定理}
\newtheorem{lemma}{引理}
\newtheorem{proposition}{命题}
\newtheorem{corollary}{推论}
\theoremstyle{definition}
\newtheorem{definition}{定义}
\newtheorem{assumption}{假设}
\newtheorem{example}{例}
\theoremstyle{remark}
\newtheorem{remark}[theorem]{注}

\newcommand{\R}{\mathbb{R}}
\newcommand{\Sm}{\mathbb{S}}
\newcommand{\bq}{\mathbf{q}}
\newcommand{\bt}{\mathbf{t}}
\newcommand{\bx}{\mathbf{x}}
\newcommand{\bu}{\mathbf{u}}
\newcommand{\br}{\mathbf{r}}
\newcommand{\bM}{\mathbf{M}}
\newcommand{\bD}{\mathbf{D}}
\newcommand{\bG}{\mathbf{G}}
\newcommand{\bH}{\mathbf{H}}
\newcommand{\bE}{\mathbf{E}}
\newcommand{\bJ}{\mathbf{J}}
\newcommand{\vtwo}{v_2}
\newcommand{\tr}{\operatorname{tr}}
\newcommand{\rank}{\operatorname{rank}}
\newcommand{\diag}{\operatorname{diag}}
\newcommand{\Proj}{\operatorname{Proj}}
\newcommand{\argmin}{\operatorname*{arg\,min}}

\newcommand{\pend}[1]{\textcolor{red}{[#1]}}
\newcommand{\ev}[1]{\emph{#1}}
.
\usepackage{amsmath,amssymb,amsthm,mathtools}
\usepackage{graphicx}
\usepackage{booktabs,multirow,array,tabularx}
\usepackage[caption=false]{subfig}
\usepackage{algorithm}
\usepackage{algpseudocode}
\usepackage[margin=2cm]{geometry}
\usepackage{url}
\usepackage{textcomp}
\usepackage{microtype}
\usepackage{xcolor}
\usepackage{bm}
\usepackage[hidelinks]{hyperref}

\setCJKmainfont{Noto Serif CJK SC}
\setCJKsansfont{Noto Sans CJK SC}
\setCJKmonofont{Noto Sans CJK SC}

\graphicspath{{figures/}}
\linespread{1.15}
\raggedbottom
\setlength{\emergencystretch}{2em}
\setlength{\columnsep}{0.75cm}
\renewcommand{\figurename}{图}
\renewcommand{\tablename}{表}
\renewcommand{\refname}{参考文献}
\floatname{algorithm}{算法}
\algrenewcommand\algorithmicrequire{\textbf{输入：}}
\algrenewcommand\algorithmicensure{\textbf{输出：}}

% Independent counters for lemmas / theorems / propositions so that the numbering
% matches the English source (Lemma 1, Theorem 1, Lemmas 2--5, Theorems 2--3,
% Propositions 1--3), see en/preamble.tex.
\newtheorem{theorem}{定理}
\newtheorem{lemma}{引理}
\newtheorem{proposition}{命题}
\newtheorem{corollary}{推论}
\theoremstyle{definition}
\newtheorem{definition}{定义}
\newtheorem{assumption}{假设}
\newtheorem{example}{例}
\theoremstyle{remark}
\newtheorem{remark}[theorem]{注}

\newcommand{\R}{\mathbb{R}}
\newcommand{\Sm}{\mathbb{S}}
\newcommand{\bq}{\mathbf{q}}
\newcommand{\bt}{\mathbf{t}}
\newcommand{\bx}{\mathbf{x}}
\newcommand{\bu}{\mathbf{u}}
\newcommand{\br}{\mathbf{r}}
\newcommand{\bM}{\mathbf{M}}
\newcommand{\bD}{\mathbf{D}}
\newcommand{\bG}{\mathbf{G}}
\newcommand{\bH}{\mathbf{H}}
\newcommand{\bE}{\mathbf{E}}
\newcommand{\bJ}{\mathbf{J}}
\newcommand{\vtwo}{v_2}
\newcommand{\tr}{\operatorname{tr}}
\newcommand{\rank}{\operatorname{rank}}
\newcommand{\diag}{\operatorname{diag}}
\newcommand{\Proj}{\operatorname{Proj}}
\newcommand{\argmin}{\operatorname*{arg\,min}}

\newcommand{\pend}[1]{\textcolor{red}{[#1]}}
\newcommand{\ev}[1]{\emph{#1}}
.
\usepackage{amsmath,amssymb,amsthm,mathtools}
\usepackage{graphicx}
\usepackage{booktabs,multirow,array,tabularx}
\usepackage[caption=false]{subfig}
\usepackage{algorithm}
\usepackage{algpseudocode}
\usepackage[margin=2cm]{geometry}
\usepackage{url}
\usepackage{textcomp}
\usepackage{microtype}
\usepackage{xcolor}
\usepackage{bm}
\usepackage[hidelinks]{hyperref}

\setCJKmainfont{Noto Serif CJK SC}
\setCJKsansfont{Noto Sans CJK SC}
\setCJKmonofont{Noto Sans CJK SC}

\graphicspath{{figures/}}
\linespread{1.15}
\raggedbottom
\setlength{\emergencystretch}{2em}
\setlength{\columnsep}{0.75cm}
\renewcommand{\figurename}{图}
\renewcommand{\tablename}{表}
\renewcommand{\refname}{参考文献}
\floatname{algorithm}{算法}
\algrenewcommand\algorithmicrequire{\textbf{输入：}}
\algrenewcommand\algorithmicensure{\textbf{输出：}}

% Independent counters for lemmas / theorems / propositions so that the numbering
% matches the English source (Lemma 1, Theorem 1, Lemmas 2--5, Theorems 2--3,
% Propositions 1--3), see en/preamble.tex.
\newtheorem{theorem}{定理}
\newtheorem{lemma}{引理}
\newtheorem{proposition}{命题}
\newtheorem{corollary}{推论}
\theoremstyle{definition}
\newtheorem{definition}{定义}
\newtheorem{assumption}{假设}
\newtheorem{example}{例}
\theoremstyle{remark}
\newtheorem{remark}[theorem]{注}

\newcommand{\R}{\mathbb{R}}
\newcommand{\Sm}{\mathbb{S}}
\newcommand{\bq}{\mathbf{q}}
\newcommand{\bt}{\mathbf{t}}
\newcommand{\bx}{\mathbf{x}}
\newcommand{\bu}{\mathbf{u}}
\newcommand{\br}{\mathbf{r}}
\newcommand{\bM}{\mathbf{M}}
\newcommand{\bD}{\mathbf{D}}
\newcommand{\bG}{\mathbf{G}}
\newcommand{\bH}{\mathbf{H}}
\newcommand{\bE}{\mathbf{E}}
\newcommand{\bJ}{\mathbf{J}}
\newcommand{\vtwo}{v_2}
\newcommand{\tr}{\operatorname{tr}}
\newcommand{\rank}{\operatorname{rank}}
\newcommand{\diag}{\operatorname{diag}}
\newcommand{\Proj}{\operatorname{Proj}}
\newcommand{\argmin}{\operatorname*{arg\,min}}

\newcommand{\pend}[1]{\textcolor{red}{[#1]}}
\newcommand{\ev}[1]{\emph{#1}}
